\documentclass[11pt]{article}
\usepackage[a4paper,margin=1in]{geometry}
\usepackage{amsmath,amssymb,mathtools,booktabs,enumitem}
\usepackage[dvipsnames]{xcolor}
\usepackage{hyperref,microtype,comment}
\hypersetup{colorlinks=true,linkcolor=MidnightBlue,urlcolor=MidnightBlue,
 pdftitle={Recitation 9: Two-Sample Tests},pdfauthor={Xuanxi Zhang}}
% Default: student handout. Define \WITHSOLUTIONS for answers and pacing.
\newif\ifsolutions
\ifdefined\WITHSOLUTIONS\solutionstrue\else\solutionsfalse\fi
\newcommand{\rv}[1]{\widetilde{#1}}
\newcommand{\E}{\mathbb E}
\newcommand{\Pbb}{\mathbb P}
\newcommand{\Var}{\operatorname{Var}}
\newcommand{\se}{\operatorname{se}}
\ifsolutions
 \newenvironment{answer}{\par\medskip\noindent\textcolor{MidnightBlue}{\textbf{Solution.}}\ }{\par\medskip}
 \newcommand{\workingspace}[1]{}
\else
 \excludecomment{answer}
 \newcommand{\workingspace}[1]{\vspace{#1}}
\fi
\setlist[enumerate]{label=(\alph*),leftmargin=2em,itemsep=0.4em}
\setlist[itemize]{leftmargin=1.5em,itemsep=0.3em}
\title{\huge Recitation 9\\[0.35em]\normalsize Two-Sample Tests\\[0.35em]
 \large\ifsolutions Solutions and Teaching Notes\else Review and Practice Problems\fi}
\author{Xuanxi Zhang}
\date{November 6, 2026}
\begin{document}
\maketitle
\noindent\textit{Goal:} Compare an observed difference between two groups with
the variation we would expect if their population parameters were equal.

\section*{Review: Antetokounmpo's free throws}
\paragraph{1. Turn the question into a model.}
The \href{https://github.com/cfgranda/ps4ds/blob/main/slides/hypothesis\%20testing/two_sample_tests_handout.pdf}{lecture handout} reports the following free throws:
\begin{center}
\begin{tabular}{@{}lrrr@{}}
\toprule Location & Made & Missed & Attempts\\\midrule
Home ($H$) & 34 & 10 & 44\\
Away ($A$) & 22 & 19 & 41\\\bottomrule
\end{tabular}
\end{center}
Let $\rv{x}_{H,i}\sim\operatorname{Bernoulli}(\theta_H)$ and
$\rv{x}_{A,j}\sim\operatorname{Bernoulli}(\theta_A)$ indicate made shots.
Assume independent shots, with a constant success probability within each
location. The parameters $\theta_H,\theta_A$ are fixed but unknown.
For ``Is his free-throw percentage \emph{different} at home and away?'', use
\[
 H_0:\theta_H=\theta_A=\theta,\qquad
 H_1:\theta_H\ne\theta_A.
\]
This is a \textbf{two-sample} problem: both success probabilities are unknown
and both are estimated from data. For Bernoulli data, equal means also imply
equal distributions; this is not true for arbitrary distributions.

\paragraph{2. Measure the observed difference.}
Let $X_H$ and $X_A$ denote the observed home and away datasets. Then
\[
 m(X_H)=\frac{34}{44}\approx0.773,\qquad
 m(X_A)=\frac{22}{41}\approx0.537,
\]
\[
 d_{\mathrm{obs}}=m(X_H)-m(X_A)\approx0.236.
\]
The gap is \textbf{23.6 percentage points}. Is it large relative to sampling
variation? The counts alone do not answer that question.

\paragraph{3. Work out the variation under the null.}
Before observing the data, let
$\rv m_H=44^{-1}\sum_i\rv x_{H,i}$,
$\rv m_A=41^{-1}\sum_j\rv x_{A,j}$, and $\rv d=\rv m_H-\rv m_A$.
Under $H_0$, independence gives
\[
 \E_0[\rv d]=0,\qquad
 \Var_0(\rv d)=\theta(1-\theta)
 \left(\frac1{44}+\frac1{41}\right).
\]
\textbf{The variances add even though the means are subtracted.}
The standard error is the square root of this variance.

\newpage
\section*{Review: from a difference to a decision}
\paragraph{4. Estimate the null standard error and standardize.}
Under $H_0$, all 85 attempts share one success probability. Estimate it by
pooling the successes and attempts:
\[
 \theta_{\mathrm{ML},0}=\frac{34+22}{44+41}=\frac{56}{85}\approx0.659,
 \qquad
 \se_{0,\mathrm{est}}=
 \sqrt{\theta_{\mathrm{ML},0}(1-\theta_{\mathrm{ML},0})
 \left(\frac1{44}+\frac1{41}\right)}\approx0.1029.
\]
Thus $z_{\mathrm{obs}}=d_{\mathrm{obs}}/\se_{0,\mathrm{est}}\approx2.295$:
the observed gap is about 2.3 estimated null standard errors above zero.
For sufficiently large samples, the corresponding random standardized
statistic $\rv z_{\mathrm{test}}$ is approximately $\mathcal N(0,1)$ under $H_0$.

\paragraph{5. Choose the tail that matches the question.}
Write $\Phi(u)=\Pbb(\rv z\leq u)$ for $\rv z\sim\mathcal N(0,1)$.
Choose the alternative and significance level \emph{before} inspecting outcomes.
\begin{center}
\begin{tabular}{@{}lll@{}}
\toprule Alternative & Approximate $p$ value & Reject at $\alpha=0.05$ if\\\midrule
$\theta_H>\theta_A$ & $1-\Phi(z_{\mathrm{obs}})$ & $z_{\mathrm{obs}}\geq1.645$\\
$\theta_H<\theta_A$ & $\Phi(z_{\mathrm{obs}})$ & $z_{\mathrm{obs}}\leq-1.645$\\
$\theta_H\ne\theta_A$ & $2[1-\Phi(|z_{\mathrm{obs}}|)]$ & $|z_{\mathrm{obs}}|\geq1.960$\\
\bottomrule
\end{tabular}
\end{center}
For the original question, both directions count as evidence:
\[
 p_{\mathrm{two}}\approx\Pbb(|\rv z|\geq2.295)\approx0.0218<0.05.
\]
Reject $H_0$ at the 5\% level. A \emph{prespecified} home-is-better alternative
would instead give $p_{\mathrm{right}}\approx0.0109$.

\paragraph{6. Say what the result means.}
Under the equal-probability model, a standardized gap at least this extreme
in either direction has approximate probability 2.2\%. This is evidence of
a difference; it is \emph{not} a 2.2\% probability that $H_0$ is true.
It does not show that crowd taunts caused the difference: location was not
randomly assigned, and opponents, fatigue, and game conditions may differ.

\paragraph{Check the approximation.}
A common rule of thumb asks for at least 10 expected makes and 10 expected
misses in each group, using the pooled estimate. Here these are
$(28.99,15.01)$ at home and $(27.01,13.99)$ away. This supports a normal
approximation, but does not establish independence. Shots within a game may
be dependent. The nominal false-positive rate is \emph{approximate} for this
plug-in $z$ test and relies on the model.

\paragraph{The same idea for other means.}
For two independent samples of sizes $n_A,n_B$, with population means
$\mu_A,\mu_B$ and \emph{known} variances $\sigma_A^2,\sigma_B^2$,
\[
 \se(\rv m_B-\rv m_A)=\sqrt{\frac{\sigma_B^2}{n_B}+\frac{\sigma_A^2}{n_A}},
 \qquad z_{\mathrm{obs}}=
 \frac{m(X_B)-m(X_A)}{\sqrt{\sigma_B^2/n_B+\sigma_A^2/n_A}}
 \quad(H_0:\mu_B=\mu_A).
\]
This is an exact Gaussian test for independent Gaussian observations with
known variances; otherwise a CLT may justify an approximation. With unknown
variances and small samples, a suitable two-sample $t$ procedure is needed.

\newpage
\section*{Practice problems}
\noindent Three core problems; Problem 4 is a backup. A calculator is useful.
Use $\Phi(1.645)\approx0.95$, $\Phi(1.960)\approx0.975$, and
$\Phi(2.295)\approx0.9891$.

\subsection*{Problem 1: Why do we pool the free throws?}
\noindent\textit{Independent work: 8 minutes. Based on the lecture example.}
Use the home/away data on page 1 and the independent Bernoulli model.
\begin{enumerate}
 \item Write the likelihood under $H_0:\theta_H=\theta_A=\theta$, up to a
 factor independent of $\theta$. Derive $\theta_{\mathrm{ML},0}$.
 Why not use the unweighted average of the two observed success proportions?
 \item Starting from $\Var(\rv x_{H,i})=\theta(1-\theta)$ under $H_0$,
 derive $\Var_0(\rv d)$. Explain both denominators and the plus sign.
 \item Compute the estimated null SE, standardized statistic, and two-sided
 $p$ value. State the decisions at $\alpha=0.05$ and $\alpha=0.01$.
 \item A student treats $22/41$ as the \emph{known} away success probability
 and includes only the home group's variance contribution in the SE.
 What source of uncertainty is missing?
\end{enumerate}
\workingspace{9cm}
\begin{answer}
\textbf{(a)} There are 56 makes and 29 misses in total, so
\[
 L_0(\theta)\propto\theta^{56}(1-\theta)^{29},\qquad
 \ell_0'(\theta)=\frac{56}{\theta}-\frac{29}{1-\theta}.
\]
Setting the derivative to zero gives
$\boxed{\theta_{\mathrm{ML},0}=56/85}$.
The second derivative is negative on $(0,1)$, and the likelihood vanishes
at both endpoints, so this is the maximum. Equivalently,
\[
 \theta_{\mathrm{ML},0}
 =\frac{44\,m(X_H)+41\,m(X_A)}{85}.
\]
Each attempt receives the same weight. An unweighted average would give
equal weight to groups with different numbers of attempts.

\textbf{(b)} Independence within each group gives
$\Var_0(\rv m_H)=44\theta(1-\theta)/44^2=\theta(1-\theta)/44$;
similarly, $\Var_0(\rv m_A)=\theta(1-\theta)/41$.
Between-group independence makes their covariance zero, hence
\[
 \boxed{\Var_0(\rv d)=\Var_0(\rv m_H)+\Var_0(\rv m_A)
 =\theta(1-\theta)(1/44+1/41).}
\]
The minus sign is squared when computing the variance.

\textbf{(c)} Using unrounded fractions,
\[
 d_{\mathrm{obs}}=0.23614,\quad \se_{0,\mathrm{est}}=0.10291,\quad
 z_{\mathrm{obs}}=2.29460,\quad p_{\mathrm{two}}\approx0.02176.
\]
Reject at 5\%; do not reject at 1\%. Failure to reject at 1\% does not
establish equal success probabilities.

\textbf{(d)} The away proportion is itself estimated from only 41 attempts.
Its variance contribution $\theta(1-\theta)/41$ is missing. Using the same
pooled plug-in value, the proposed SE is about $0.07147$, too small; it
overstates the evidence against $H_0$.
\end{answer}

\newpage
\subsection*{Problem 2: Comparing treatment and control means}
\noindent\textit{Independent work: 8 minutes. Archived Exercise 11.19.}
An experiment compares a control group ($A$) with an independent treatment
group ($B$). Observations are approximately normally distributed, and the
standard deviation of an individual measurement in either group is known
to be $\sigma=5$. Each group has 25 observations.
\begin{enumerate}
 \item What is the standard error of $\rv m_B-\rv m_A$?
 \item For $H_0:\mu_B=\mu_A$ versus $H_1:\mu_B>\mu_A$, give the
 level-$0.05$ rejection region, both in standardized units and in the
 original measurement units.
 \item Suppose the test's $p$ value is $0.07$. Would you reject at
 $\alpha=0.10$? At $\alpha=0.05$?
 \item What is the level-$0.05$ rejection region if instead
 $H_1:\mu_B\ne\mu_A$? Explain why the cutoff changes.
\end{enumerate}
\noindent\textit{Discussion:} Why is $5/\sqrt{50}$ not the SE in part (a)?
\workingspace{10cm}
\begin{answer}
\textbf{(a)} Each sample mean has variance $5^2/25=1$. Independence gives
\[
 \boxed{\se(\rv m_B-\rv m_A)=\sqrt{\frac{25}{25}+\frac{25}{25}}
 =\sqrt2\approx1.414.}
\]
The value $5/\sqrt{50}$ is the SE of the grand mean of 50 independent
observations with variance 25, not the difference of two group means.

\textbf{(b)} Under $H_0$, the standardized difference is approximately
standard Gaussian. Let $d_{\mathrm{obs}}=m(X_B)-m(X_A)$. Reject when
\[
 \boxed{z_{\mathrm{obs}}=\frac{d_{\mathrm{obs}}}{\sqrt2}\geq1.645}
 \qquad\Longleftrightarrow\qquad
 \boxed{d_{\mathrm{obs}}\geq1.645\sqrt2\approx2.326.}
\]
Only a sufficiently positive difference supports this alternative.

\textbf{(c)} Since $0.07\leq0.10$, reject at 10\%. Since $0.07>0.05$,
do not reject at 5\%. Choose the significance level in advance;
changing it after seeing the $p$ value changes the planned decision rule.

\textbf{(d)} Split the 5\% null tail area equally between the two directions:
\[
 \boxed{|z_{\mathrm{obs}}|\geq1.960}
 \qquad\Longleftrightarrow\qquad
 \boxed{|d_{\mathrm{obs}}|\geq1.960\sqrt2\approx2.772.}
\]
Thus reject if $d_{\mathrm{obs}}\leq-2.772$ or $d_{\mathrm{obs}}\geq2.772$.
The positive cutoff is larger because only 2.5\% of the null distribution
lies in the upper rejection tail.

\medskip
\noindent\textit{Connection to the free throws:} Both tests divide a difference
by its null standard error. Here the variances are known. In the Bernoulli
example, the unknown common success probability determines both variances
and is estimated using pooled data.
\end{answer}

\newpage
\subsection*{Problem 3: What does the test tell us?}
\noindent\textit{Independent work: 7 minutes. Adapted from archived Exercise
9.26(b)--(d), with extensions to the free-throw example.}
Mark each statement true or false and justify your answer in one or two
sentences. For a false statement, give a correct replacement.
\begin{enumerate}
 \item If the $p$ value is $0.03$, the test rejects at significance level $0.02$.
 \item If a test using the rule ``reject when $p\leq\alpha$'' rejects at
 significance level $0.06$, then its $p$ value is at most $0.06$.
 \item The free-throw test's two-sided $p$ value $0.0218$ is the probability
 that the home and away success probabilities are equal.
 \item If we compute away minus home instead of home minus away, the
 two-sided $p$ value is unchanged.
 \item After seeing that the home percentage is higher, we may replace
 the planned two-sided test by a one-sided test and halve its $p$ value
 while keeping the advertised significance level.
 \item Rejecting equality establishes that crowd taunts caused the lower
 away success rate.
\end{enumerate}
\workingspace{9cm}
\begin{answer}
\textbf{(a) False.} $0.03>0.02$, so do not reject. The same result would
reject at a prespecified 5\% level.

\textbf{(b) True.} This is exactly the stated rejection rule. A smaller
significance level might lead to a different decision.

\textbf{(c) False.} The $p$ value measures how extreme the data are under
the null model. Here it is approximately the probability, under $H_0$, of
a standardized difference with magnitude at least $2.295$. It is not a
posterior probability of $H_0$.

\textbf{(d) True.} Reversing the subtraction changes $z_{\mathrm{obs}}$
to $-z_{\mathrm{obs}}$ but leaves $|z_{\mathrm{obs}}|$ unchanged.
For a one-sided test, the sign and the alternative must be kept consistent:
``home is better'' becomes a \emph{left}-tail test for away minus home.

\textbf{(e) False.} The direction must be chosen before examining the
outcomes. If we always choose whichever direction the observed effect
favors and then use a one-sided 5\% cutoff, we reject for
$|z_{\mathrm{obs}}|\geq1.645$. Under an ideal continuous standard-normal
null, that event has probability 10\%, not 5\%.

\textbf{(f) False.} A difference in probabilities is not a causal
attribution. This observational comparison does not isolate crowd taunts
from other differences between home and away games. The independent-shot
model may also be imperfect.

\medskip
\noindent\textit{Exit check:} Complete the sentence: ``Assuming the null
model, the $p$ value is the probability of \underline{\hspace{3cm}}.''
The missing phrase is ``a test statistic at least as extreme as the one
observed, in the direction(s) specified by the alternative.''
\end{answer}

\newpage
\subsection*{Problem 4 (backup): Is the normal approximation credible?}
\noindent\textit{Independent work: 5 minutes. Adapted from archived
Recitation Worksheet 11, Problem 2.}
A study records exposure to an environmental factor and disease status:
\begin{center}
\begin{tabular}{@{}lrr@{}}
\toprule & Not exposed & Exposed\\\midrule
No disease & 128 & 7\\
Disease & 19 & 7\\\bottomrule
\end{tabular}
\end{center}
Model individuals as independent, with disease probabilities
$\theta_N$ and $\theta_E$ in the two groups.
\begin{enumerate}
 \item State hypotheses for whether the disease probabilities differ, and
 find the observed exposed-minus-unexposed difference.
 \item Find the pooled null estimate and the expected disease/no-disease
 counts in each group using that estimate.
 \item Does the ``at least 10 expected successes and failures per group''
 guideline hold? Explain why computing a normal-tail $p$ value alone
 would not settle whether the approximation is reliable.
\end{enumerate}
\workingspace{9cm}
\begin{answer}
\textbf{(a)} Test $H_0:\theta_E=\theta_N$ against
$H_1:\theta_E\ne\theta_N$. The observed difference is
\[
 \boxed{d_{\mathrm{obs}}=\frac7{14}-\frac{19}{147}\approx0.37075.}
\]
The difference is about 37.1 percentage points.

\textbf{(b)} The common-probability estimate is
$\theta_{\mathrm{ML},0}=26/161\approx0.16149$.
Multiplying each group size by this value and its complement gives
\begin{center}
\begin{tabular}{@{}lrr@{}}
\toprule Expected count under fitted $H_0$ & Not exposed & Exposed\\\midrule
Disease & 23.74 & 2.26\\
No disease & 123.26 & 11.74\\\bottomrule
\end{tabular}
\end{center}
\textbf{(c)} No: the expected exposed disease count is only 2.26.
A mechanical pooled $z$ calculation gives $z\approx3.602$ and an approximate
two-sided $p\approx0.000316$, but that small number does not validate its
Gaussian calibration. The guideline is a diagnostic, not a theorem
asserting failure or success. A suitable exact method, such as Fisher's
exact test for this $2\times2$ table, avoids this normal approximation;
its calculation is beyond this recitation.

Even a well-calibrated test would assess association, not establish that
exposure caused disease. The original archive asks for a significance
test; this adaptation emphasizes checking the assumptions first.
\end{answer}

\ifsolutions
\newpage
\section*{Instructor notes: 75-minute plan}
\begin{center}
\begin{tabular}{@{}p{0.16\textwidth}p{0.77\textwidth}@{}}
\toprule Time & Activity\\\midrule
0--7 min & Free-throw table; two unknown probabilities; formulate a
two-sided alternative before calculating.\\
7--16 min & Derive the difference's null mean and variance; explain
pooling and standardization.\\
16--25 min & One versus two tails; $p$ value, decision, approximation
checks, and the limits of a causal conclusion.\\
25--42 min & Problem 1: 8 minutes independent work, 9 minutes discussion.\\
42--58 min & Problem 2: 8 minutes independent work, 8 minutes discussion.\\
58--72 min & Problem 3: 7 minutes independent work, 7 minutes discussion.\\
72--75 min & Exit check and questions.\\\bottomrule
\end{tabular}
\end{center}
\noindent\textit{Fast-class option:} Keep all independent-work periods.
Shorten discussions by 2 minutes each for Problems 1 and 2 and by 1 minute
for Problem 3. Then use minutes 67--72 for Problem 4 and 72--75 for its
discussion and the exit check.

\paragraph{Board sequence.}
Write the two proportions, then ask whether a 23.6-point gap is surprising.
Build the chain
\[
 \text{equal probabilities}\ \Longrightarrow\ \E_0[\rv d]=0
 \ \Longrightarrow\ \se_{0,\mathrm{est}}\approx0.103
 \ \Longrightarrow\ z\approx2.295
 \ \Longrightarrow\ p_{\mathrm{two}}\approx0.022.
\]
Sketch a bell curve with both tails beyond $\pm2.295$ shaded. Ask which
tail would count for a prespecified ``home is better'' question. Do not
describe the positive raw difference alone as the two-sided evidence:
its \emph{magnitude} is what matters.

\paragraph{Assumptions and notation to emphasize.}
The true common parameter $\theta$ is deterministic and unknown;
$\theta_{\mathrm{ML},0}=56/85$ is its observed estimate under $H_0$.
The random means have tildes, and $X_H,X_A$ denote observed datasets.
The pooled null SE is appropriate for testing equal Bernoulli probabilities.
An unrestricted confidence interval for their difference commonly uses
separate group estimates instead; do not conflate the two calculations.
Independent groups do not by themselves ensure independence of shots
\emph{within} groups. A game-level dependence model would require more
information than these two totals.

\paragraph{Sources and adaptations.}
The review and Problem 1 follow Carlos Fernandez-Granda's
\href{https://github.com/cfgranda/ps4ds/blob/main/slides/hypothesis\%20testing/two_sample_tests_handout.pdf}{\emph{Two-Sample Tests} handout},
especially the free-throw counts and the two tail probabilities.
We call the square-root quantity a standard error (the slide labels it
``variance'') and describe Gaussian calibration as approximate.
Problem 2 preserves Exercise 11.19 from the Fall 2025 archive's
\texttt{mine/r10/main.tex}; it also appears in archived Worksheet 9.
Problem 3(a)--(c) adapts Exercise 9.26(b)--(d) from
\texttt{mine/r8/main.tex}; (d)--(f) connect these ideas to the lecture example.
Problem 4 preserves the counts in archived Worksheet 11, Problem 2,
and changes the task to checking normal-approximation conditions.
The archived worksheets themselves carry Spring 2025 headers.
\fi
\end{document}
